% Generated by scripts/make_paper_tables.py from the frozen results.
% Do not edit by hand; the next build overwrites it.
\begin{table}[tbp]
\centering
\caption{Candidate explanations for the gap, one input at a time, per \$1,000 of principal}
\label{tab:inputs}
\small
\begin{tabular}{l r r r r r}
\toprule
Note & Gap to close & +0.05 correlation & Flat vol & Timing & Call cost \\
\midrule
N1 & $-$27.92 & 9.14 & 62.49 & $-$4.04 & 27.99 \\
N2 & $-$48.72 & 7.70 & 70.21 & $-$7.94 & 22.34 \\
N3 & $-$20.84 & 5.58 & 44.50 & $-$4.56 & 40.09 \\
N4 & $-$19.16 & 3.55 & 31.90 & $-$3.64 & 36.38 \\
N5 & $-$11.41 & 3.32 & 6.65 & $-$3.76 & $-$51.99 \\
N6 & $-$13.29 & 3.10 & 26.71 & $-$4.55 & 34.23 \\
\bottomrule
\end{tabular}
\notetext{\textit{Notes:} Each column is the change in model value from moving one input while the rest are held at their calibrated values, on a common set of paths. \textit{+0.05 correlation} raises every pairwise correlation, confirming the investor is long correlation on every note. \textit{Flat vol} replaces the fitted surface with a single at-the-money volatility per leg. It raises every value, and by more than the gap on five of the six, so the shape of the surface is the one input measured here whose effect both runs toward the disclosed value and is large enough to span the difference. \textit{Timing} aligns the surface to the pricing date rather than the snapshot date and lowers every value, widening the gap. \textit{Call cost} is what the issuer's option to redeem is worth to the issuer, from the decomposition in Table~\ref{tab:decomposition}. No single column closes the gap in the first; \S\ref{sec:inputs} takes them in turn.}
\end{table}
